% Meio-circuito em modo comum de M2 (Questão 3), com R_S representado por 2R_S, e seu modelo pi com r_o.
% Compilar com pdflatex; o PDF resultante vai para figuras/fig_q3_cm_meio.pdf.
\documentclass[11pt,border=4pt]{standalone}
% Estilo comum dos esquemáticos do relatório (circuitikz).
% Uso: \documentclass[11pt,border=4pt]{standalone} e, em seguida, % Estilo comum dos esquemáticos do relatório (circuitikz).
% Uso: \documentclass[11pt,border=4pt]{standalone} e, em seguida, % Estilo comum dos esquemáticos do relatório (circuitikz).
% Uso: \documentclass[11pt,border=4pt]{standalone} e, em seguida, \input{estilo_circuitos}.
\usepackage[T1]{fontenc}
\usepackage[utf8]{inputenc}
\usepackage[brazilian]{babel}
\usepackage{lmodern}
\usepackage{amsmath}
\usepackage[american]{circuitikz}
\usetikzlibrary{calc,arrows.meta}

% Cores das figuras originais
\definecolor{anot}{RGB}{0,135,200}      % anotações: valores calculados, setas, resistências vistas
\definecolor{sombra}{RGB}{220,242,253}  % região do transistor nos modelos T e pi
\definecolor{term}{RGB}{62,124,148}     % nomes de terminais destacados (B, C, E)

% Fios com 0,4 pt e corpo dos componentes com 0,8 pt, como nas figuras originais
\ctikzset{bipoles/thickness=2}
.
\usepackage[T1]{fontenc}
\usepackage[utf8]{inputenc}
\usepackage[brazilian]{babel}
\usepackage{lmodern}
\usepackage{amsmath}
\usepackage[american]{circuitikz}
\usetikzlibrary{calc,arrows.meta}

% Cores das figuras originais
\definecolor{anot}{RGB}{0,135,200}      % anotações: valores calculados, setas, resistências vistas
\definecolor{sombra}{RGB}{220,242,253}  % região do transistor nos modelos T e pi
\definecolor{term}{RGB}{62,124,148}     % nomes de terminais destacados (B, C, E)

% Fios com 0,4 pt e corpo dos componentes com 0,8 pt, como nas figuras originais
\ctikzset{bipoles/thickness=2}
.
\usepackage[T1]{fontenc}
\usepackage[utf8]{inputenc}
\usepackage[brazilian]{babel}
\usepackage{lmodern}
\usepackage{amsmath}
\usepackage[american]{circuitikz}
\usetikzlibrary{calc,arrows.meta}

% Cores das figuras originais
\definecolor{anot}{RGB}{0,135,200}      % anotações: valores calculados, setas, resistências vistas
\definecolor{sombra}{RGB}{220,242,253}  % região do transistor nos modelos T e pi
\definecolor{term}{RGB}{62,124,148}     % nomes de terminais destacados (B, C, E)

% Fios com 0,4 pt e corpo dos componentes com 0,8 pt, como nas figuras originais
\ctikzset{bipoles/thickness=2}

\begin{document}
\begin{circuitikz}
  % meia largura do zigue-zague do resistor, raio dos terminais e folga padrão entre rótulo e traço
  \def\zw{0.21cm}
  \def\rc{1.6pt}
  \def\gap{4.5pt}
  % ---------------- (a) meio-circuito em modo comum (nó do dreno em (0,0)) ----------------
  \draw (0,1.6) node[ground, rotate=180]{} to[R] (0,0) -- (0,-0.4)
        node[nmos, arrowmos, anchor=D] (M) {};
  \draw (0,0) node[circ]{} to[short, -o] (1.2,0);
  \draw (M.G) -- ++(-0.7,0) node[ocirc] (VIa) {};
  \draw (M.S) to[R] ++(0,-1.8) node[ground]{};
  \draw[anot, line width=0.8pt, -latex] ($(M.S)+(0.45,-0.45)$) -- ++(0,-0.7);
  % rótulos de (a)
  \node[anchor=east, inner sep=0pt] at ([xshift=-\zw-\gap]0,0.8) {$R_D$};
  \node[anot, anchor=west, inner sep=0pt] at ([xshift=\rc+\gap]1.2,0) {$v_{o2}$};
  \node[anot, anchor=east, inner sep=0pt] at ([xshift=-\rc-\gap]VIa.center) {$v_{icm}$};
  \node[anchor=west, inner sep=0pt] at ([xshift=0.27cm]M.G -| M.D) {$M_2$};
  \node[anchor=east, inner sep=0pt] at ([xshift=-\zw-\gap]$(M.S)+(0,-0.9)$) {$2R_S$};
  \node[anot, anchor=west, inner sep=0pt] at ([xshift=6pt]$(M.S)+(0.45,-0.8)$) {$i$};

  % ---------------- (b) modelo pi (nó D em (7.6,0)) ----------------
  \begin{scope}[shift={(7.6,0)}, local bounding box=bb]
    \draw (0,0) to[cI] (0,-2.4);
    \draw (0,0) -- (1.4,0) to[R] (1.4,-2.4) -- (0,-2.4);
    \draw (1.4,0) -- (2.7,0) to[short, -o] (3.6,0);
    \draw (2.7,0) to[R] (2.7,1.7) node[ground, rotate=180]{};
    % fio entre S e 2R_S mais longo que na original, para o rótulo do nó S caber entre o fio e o resistor
    \draw (0.7,-2.4) -- (0.7,-2.6) to[R] (0.7,-4.6) node[ground]{};
    \draw (0,0) node[circ]{} (1.4,0) node[circ]{} (2.7,0) node[circ]{} (0.7,-2.4) node[circ]{};
    \draw (-2.3,-1.2) node[ocirc] (G) {};
    % rótulos de (b)
    \node[anchor=south, inner sep=0pt] at ([yshift=\rc+6pt]G.center) {$G$};
    \node[anot, anchor=north, inner sep=0pt] at ([yshift=-\rc-7pt]G.center) {$v_{icm}$};
    \node[anchor=south, inner sep=0pt] at ([yshift=\rc+6pt]0,0) {$D$};
    \node[anchor=east, inner sep=0pt] at ([xshift=-0.49cm-\gap]0,-1.2) {$g_mv_{gs}$};
    \node[anchor=west, inner sep=0pt] at ([xshift=\zw+\gap]1.4,-1.2) {$r_o$};
    \node[anchor=north west, inner sep=0pt] at ([xshift=\gap, yshift=-\gap]0.7,-2.4) {$S\ (v_s)$};
    \node[anchor=east, inner sep=0pt] at ([xshift=-\zw-\gap]0.7,-3.6) {$2R_S$};
    \node[anchor=west, inner sep=0pt] at ([xshift=\zw+\gap]2.7,0.85) {$R_D$};
    \node[anot, anchor=west, inner sep=0pt] at ([xshift=\rc+\gap]3.6,0) {$v_{o2}$};
    \coordinate (terra) at (0.7,-4.6);
  \end{scope}
  % equações e legendas das partes, centradas como na original
  \node[anot, anchor=north, inner sep=0pt] (eq) at ([yshift=-0.42cm-4pt-5pt]bb.center |- terra)
    {$v_{gs} = v_{icm} - v_s,\quad v_s = 2R_S\,i_d$};
  \node[anchor=north, inner sep=0pt] (capb) at ([yshift=-5pt]eq.south) {(b) Modelo $\pi$-híbrido};
  \node[anchor=base, inner sep=0pt] at (capb.base -| 0,0) {(a) Meio-circuito em modo comum};
\end{circuitikz}
\end{document}
